
A large language model (LLM) scoring 95\% on a composite trustworthiness benchmark can fail reliability tests (e.g., hallucinating drug interactions), robustness tests (e.g., succumbing to adversarial prompts), and fairness tests (e.g., exhibiting demographic bias) at nearly identical rates across all three dimensions. This tight coupling raises a foundational question: are reliability, robustness, and fairness independent failure modes requiring separate interventions, or symptoms of shared underlying deficiencies?

Current trustworthiness evaluation practice treats these dimensions as independent properties. Frameworks such as HELM aggregate scores across specialized benchmarks---TrustfulQA for reliability, AdvBench for robustness, BOLD for fairness---without testing whether failures correlate across dimensions. This assumption of empirical orthogonality has practical consequences. If failures are independent, practitioners should diagnose and fix dimension-specific issues (e.g., calibration training for hallucinations, adversarial fine-tuning for robustness). If failures are tightly coupled, dimension-independent interventions may waste resources by addressing symptoms rather than shared root causes (e.g., poor training data diversity affecting all dimensions simultaneously).

Prior work provides indirect evidence of coupling. Alignment fine-tuning through reinforcement learning from human feedback (RLHF) has been observed to improve reliability, safety, and fairness concurrently. However, these observations are qualitative and lack quantitative characterization of correlation structure. No prior study has systematically tested whether trustworthiness dimensions exhibit moderate correlations (r > 0.3) expected from shared mechanisms, stronger coupling suggesting near-redundancy, or independence (r $\approx$ 0) supporting dimension-specific evaluation.

This study addresses the gap through large-scale correlation analysis across 20 LLMs spanning three model size strata (small <1B parameters, medium 1--10B, large >10B). We computed pairwise Spearman correlations between TrustfulQA, AdvBench, and BOLD scores to test the null hypothesis that trustworthiness failures are independent. Results reveal correlations far exceeding initial expectations: all three benchmark pairs exhibit r > 0.99 (p < $2\times 10^{-17})$, with coupling persisting across model scales (r > 0.98 within small, medium, and large strata).

Near-perfect correlation admits two competing explanations. First, reliability, robustness, and fairness---as operationalized by current benchmarks---may measure the same underlying construct (general model quality or training data diversity) rather than orthogonal failure mechanisms. This is consistent with observations that RLHF improves multiple dimensions without dimension-specific targeting. Second, distinct dimensions may exist but 3 benchmarks provide insufficient resolution to distinguish them---correlations could reflect shared confounds (e.g., all benchmarks sensitive to pre-training quality) rather than unified constructs. Correlation alone cannot distinguish these explanations without experimental manipulation.

To test whether correlations organize into distinct failure modes, we applied hierarchical clustering to the correlation matrix. Analysis revealed perfect cluster stability (100\% bootstrap consistency across 1000 iterations) but poor separation quality (silhouette score = 0.274 < 0.5 threshold). The divergence between stability and separation indicates that 3-benchmark evaluation cannot resolve the granularity required for failure mode taxonomy. This is a methodological limitation with implications for benchmark design: clustering k $\geq$ 3 failure modes requires n $\geq$ 3 benchmarks, yet our k = 2 evaluation (TrustfulQA+AdvBench versus BOLD) lacks both sample size and separation quality.

The study contributes:

\begin{enumerate}
\item \textbf{First large-scale empirical evidence} that trustworthiness dimensions exhibit near-perfect correlations (r > 0.99, p < $2\times 10^{-17})$ across 20 LLMs, far stronger than moderate effect sizes (r > 0.3) hypothesized from shared root causes. Coupling persists across model scales (r > 0.98 within size strata).
\end{enumerate}

\begin{enumerate}
\item \textbf{Methodological lesson} for benchmark design: 3-benchmark evaluation is insufficient to resolve distinct failure mode clusters despite perfect stability (100\% bootstrap consistency). Clustering quality thresholds (silhouette > 0.5) may be unattainable for high-correlation data (r > 0.95) even when structure is reproducible.
\end{enumerate}

\begin{enumerate}
\item \textbf{Competing explanations and testable predictions}: The Unified Capability Hypothesis posits that all three benchmarks measure the same construct, predicting that expansion to 10 benchmarks yields r > 0.95 for >80\% of pairs. The Insufficient Resolution Hypothesis posits that distinct dimensions exist but require broader measurement, predicting mean r < 0.85 with $\geq 30$\% of pairs exhibiting r < 0.7.
\end{enumerate}

Results challenge dimension-independent evaluation paradigms. If coupling generalizes beyond the 3-benchmark sample, practitioners may benefit from multi-dimensional interventions targeting shared root causes rather than dimension-specific fixes. However, interpretation remains uncertain pending broader measurement (5--10 benchmarks) and intervention validation.
